\documentclass{ximera}
\title{Substituion}
\begin{document}
\begin{abstract}
The substitution rule (the chain rule for integration) for indefinite and definite integrals, including changing the limits of integration. Textbook §5.5.
\end{abstract}
\maketitle

We're now going to look at the chain rule, but for integration. So suppose we have some (indefinite) integral like: $\int(3 x+2)^{7 / 9} d x$ ? And we're trying to find the antiderivative of $(3 x+2)^{7 / 9}$.

So far we haven't really learned how to tackle something like this, but we can use the ideas of the chain rule (and implicit differentiation) to solve something like this. What this method is called is ``substitution''.

What we want is an integral that is easier to handle. In our case, we want to change our integral to be something like $\int u^{7 / 9} d u$. Why? Because we already know how to do this function!

So if we set $u=3 x+2$, then we have $\int u^{7 / 9} d x$. But now we have a different problem. What do you think the problem is?
% open-ended question left as white space in the original; open work per spec 3(e)
\begin{freeResponse}\end{freeResponse}

Well, let's see if we can figure it out! We're gonna use implicit differentiation:
% work: differentiate both sides: du = 3 dx
\[
u=3 x+2 \Rightarrow d u=\answer{3}\, d x
\]

Therefore
% work: dx = du/3, so the integral becomes int u^{7/9} du/3; then int u^{7/9} du = u^{16/9}/(16/9) + c
\[
\begin{aligned}
\int(3 x+2)^{7 / 9} d x & =\answer{\int u^{7 / 9}\, \frac{d u}{3}} \\
& =\frac{1}{3} \int u^{7 / 9} d u \\
& =\answer{\frac{1}{3} \cdot \frac{u^{16 / 9}}{16 / 9}+c} \\
& =\frac{9}{3 \cdot 16} u^{16 / 9}+c \\
& =\frac{3}{16}(3 x+2)^{16 / 9}+c
\end{aligned}
\]

And done!

So what did we do?
% AUTHOR CHOICE: wording of the five summary steps chosen during conversion, change freely
\begin{enumerate}
  \item \wordChoice{\choice[correct]{Chose a substitution $u=g(x)$ for the inside function (here $u=3x+2$)}\choice{Found $du=g^{\prime}(x)\,dx$ (here $du=3\,dx$) and solved for $dx$}\choice{Rewrote the integral entirely in terms of $u$ and $du$}\choice{Integrated with respect to $u$}\choice{Substituted $u=g(x)$ back in to get an answer in terms of $x$}}
  \item \wordChoice{\choice{Chose a substitution $u=g(x)$ for the inside function (here $u=3x+2$)}\choice[correct]{Found $du=g^{\prime}(x)\,dx$ (here $du=3\,dx$) and solved for $dx$}\choice{Rewrote the integral entirely in terms of $u$ and $du$}\choice{Integrated with respect to $u$}\choice{Substituted $u=g(x)$ back in to get an answer in terms of $x$}}
  \item \wordChoice{\choice{Chose a substitution $u=g(x)$ for the inside function (here $u=3x+2$)}\choice{Found $du=g^{\prime}(x)\,dx$ (here $du=3\,dx$) and solved for $dx$}\choice[correct]{Rewrote the integral entirely in terms of $u$ and $du$}\choice{Integrated with respect to $u$}\choice{Substituted $u=g(x)$ back in to get an answer in terms of $x$}}
  \item \wordChoice{\choice{Chose a substitution $u=g(x)$ for the inside function (here $u=3x+2$)}\choice{Found $du=g^{\prime}(x)\,dx$ (here $du=3\,dx$) and solved for $dx$}\choice{Rewrote the integral entirely in terms of $u$ and $du$}\choice[correct]{Integrated with respect to $u$}\choice{Substituted $u=g(x)$ back in to get an answer in terms of $x$}}
  \item \wordChoice{\choice{Chose a substitution $u=g(x)$ for the inside function (here $u=3x+2$)}\choice{Found $du=g^{\prime}(x)\,dx$ (here $du=3\,dx$) and solved for $dx$}\choice{Rewrote the integral entirely in terms of $u$ and $du$}\choice{Integrated with respect to $u$}\choice[correct]{Substituted $u=g(x)$ back in to get an answer in terms of $x$}}
\end{enumerate}

\begin{exercise}\label{exercise:13-8}
With a partner, find the following integral: $\int x\left(7 x^{2}-8\right)^{16 / 7} d x$.

\textbf{Hint:} Don't worry about the $x$ by itself.
% work: u = 7x^2 - 8, du = 14x dx, so x dx = du/14;  (1/14) int u^{16/7} du = (1/14)(7/23) u^{23/7} + c = u^{23/7}/46 + c
\[
\int x\left(7 x^{2}-8\right)^{16 / 7} d x=\answer{\frac{1}{46}\left(7 x^{2}-8\right)^{23 / 7}+c}
\]
\end{exercise}

But recall that we're doing all of these things to find definite integrals! So what happens with that?

\begin{example}\label{example:13-9}
Find the following definite integral: $\int_{2}^{4} \frac{2 x-1}{x^{2}-x} d x$
% work: the natural substitution is u = x^2 - x, which gives du = (2x - 1) dx
\[
u=\answer{x^{2}-x} \Rightarrow d u=\answer{2 x-1}\, d x
\]
Therefore
% the definite integral was worked live in the original (white space) before the limits problem is raised; open work per spec 3(e)
\begin{freeResponse}\end{freeResponse}

But nowe we have a different problem!

\wordChoice{\choice[correct]{The limits $2$ and $4$ are values of $x$, not values of $u$.}\choice{The integrand is not continuous on $[2,4]$.}\choice{We forgot the constant $c$.}\choice{The substitution $u=x^{2}-x$ is not differentiable.}}

Let's look at the two functions together. The one on top (blue) is the one with variable $x$ and the one on the bottom (red) is the one with variable $u$.
\begin{image}
\includegraphics[width=0.6\textwidth]{b48a5324-93dc-45ee-b7c2-dc0b5706bda6-147_567_773_1547_666}
\end{image}

So how do we solve this? We do exactly like we did with $d x$.
% work: u = x^2 - x sends x=2 to u=2 and x=4 to u=12;  int_2^12 du/u = ln(12) - ln(2) = ln(6)
\[
x=2 \Rightarrow u=\answer{2}, \qquad x=4 \Rightarrow u=\answer{12}
\]
\[
\int_{2}^{4} \frac{2 x-1}{x^{2}-x} d x=\int_{2}^{12} \frac{1}{u} d u=\answer{\ln(6)}
\]
\end{example}

Here's the substitution rule in all its glory:

\begin{theorem}[Substitution rule for a definite integral]\label{theorem:13-10}
If $g^{\prime}(x)$ is a continuous function on the interval $[a, b]$ and if $f(x)$ is a continuous function that has an antiderivative on the interval $[g(a), g(b)]$, then
\[
\int_{a}^{b} f(g(x)) g^{\prime}(x) d x=\answer{\int_{g(a)}^{g(b)} f(u) d u}
\]
\end{theorem}

Where does the $g^{\prime}(x)$ come from? From the chain rule!
\[
\begin{aligned}
u=g(x) & \Rightarrow d u=g^{\prime}(x) d x \\
& \Rightarrow \frac{d u}{g^{\prime}(x)} d x
\end{aligned}
\]

and therefore to be able to change $f(g(x))$ to $f(u)$ we must always have $g^{\prime}(x)$ there.

\begin{example}\label{example:13-11}
Let's try this with a function that's a little bit more fun $\int \cotan(a x) d x$ for $a \in \mathbb{R}$.
% work: cotan(ax) = cos(ax)/sin(ax); u = sin(ax), du = a cos(ax) dx;  (1/a) int du/u = (1/a) ln|sin(ax)| + c
\[
\int \cotan(a x) d x=\answer{\frac{1}{a} \ln(|\sin(a x)|)+c}
\]
\end{example}

\end{document}
